\section*{Hoofdstuk \ref{ch:b3:molecular-evolution} --- Moleculaire evolutie en fylogenomica}

\begin{solution}{exo:b3:molecular-evolution:1}
$k = \mu$: het neutrale \omterm{thm:b3:molecular-evolution:neutral}{substitutietempo} per plaats per generatie is
gelijk aan het mutatietempo. Per generatie ontstaan er $2N\mu$ nieuwe
neutrale mutaties en heeft elk een kans $1/2N$ om te worden vastgelegd;
een grotere populatie maakt meer mutaties en geeft elk een evenredig
kleinere kans, en die twee effecten heffen elkaar precies op.
\end{solution}

\begin{solution}{exo:b3:molecular-evolution:2}
\omterm{def:b3:genomics:comparative}{Orthologen}: genen in twee soorten die afstammen van één gen in hun
gemeenschappelijke voorouder --- het $\beta$-globine van de mens en het
$\beta$-globine van de muis. Paralogen: genen in één genoom die van een verdubbeling
afstammen --- het $\alpha$- en het $\beta$-globine van de mens, of
$\beta$ en $\gamma$. \omterm{def:b3:molecular-evolution:duplication}{Pseudogen}: een vervallen kopie die geen eiwit meer
codeert --- het gen $\psi\beta$ in het $\beta$-globinecluster van de
mens, compleet met stopcodons.
\end{solution}

\begin{solution}{exo:b3:molecular-evolution:3}
$\omega$ vergelijkt het tempo van substitutie die het aminozuur
verandert met het tempo van stille substitutie, dat de neutrale maatstaf
is. $0.02$: sterke \omterm{def:b3:molecular-evolution:dnds}{zuiverende selectie}, waarbij \qty{98}{\%} van de
veranderingen van aminozuren wordt verwijderd --- actine, \omterm{def:b3:chromatin-epigenetics:nucleosome}{histon} H3.
$1.0$: geen beperking --- een \omterm{def:b3:molecular-evolution:duplication}{pseudogen}. $2.5$: \omterm{def:b3:molecular-evolution:dnds}{positieve selectie},
waarbij veranderingen van aminozuren worden begunstigd --- de
peptidebindende groeve van MHC, de \omterm{def:b3:virology:virus}{envelop} van hiv, het lysozym van de
herkauwer.
\end{solution}

\begin{solution}{exo:b3:molecular-evolution:4}
Eiwitten die over de zoogdieren heen worden vergeleken veranderen met
ongeveer één substitutie per plaats per miljard jaar; over een genoom
van een miljard coderende plaatsen en een generatie van enkele jaren is
dat in de orde van één vastgelegde substitutie per generatie, of één
per twee jaar. De selectiekosten van Haldane laten ongeveer één
geselecteerde substitutie per driehonderd generaties toe. De tempo's
verschillen een factor honderd of meer, dus moeten vrijwel alle
vastgelegde veranderingen neutraal zijn.
\end{solution}

\begin{solution}{exo:b3:molecular-evolution:5}
Neutraal: $1/2N = 10^{-4}$. $s = 0.001$, $4Ns = 20$: $(1 - \mathrm{e}^{-0.002})/
(1 - \mathrm{e}^{-20}) = 0.002$. $s = 0.01$: $0.0198$. $s = -0.001$: $(1 -
\mathrm{e}^{0.002})/(1 - \mathrm{e}^{20}) = 0.002\,\mathrm{e}^{-20} =
4\times 10^{-12}$. Geen ervan is effectief neutraal: daarvoor zou
$|4Ns|
\lesssim 1$ nodig zijn, d.w.z.\ $|s| \lesssim 5\times 10^{-5}$; een
selectiecoëfficiënt van een tiende procent is in een populatie van
vijfduizend al beslissend.
\end{solution}

\begin{solution}{exo:b3:molecular-evolution:6}
Ongecorrigeerd: $T = 0.12/(2\times 2\times 10^{-9}) = \qty{30}{Myr}$.
Jukes--Cantor: $d = -\tfrac{3}{4}\ln(1 - 0.16) = 0.131$, $T =
\qty{33}{Myr}$. De correctie bereikt een factor twee wanneer $-\tfrac{3}{4}
\ln(1 - 4p/3) = 2p$, bij $p \approx 0.6$: dan heeft het ruwe verschil
het meeste van zijn informatie al verloren.
\end{solution}

\begin{solution}{exo:b3:molecular-evolution:7}
$p_{N} = 7/700 = 0.010$, $p_{S} = 10/200 = 0.050$, $\omega = 0.2$:
\omterm{def:b3:molecular-evolution:dnds}{zuiverende selectie} die vier vijfde van de veranderingen van aminozuren
verwijdert. Met $35$ en $10$: $p_{N} = 0.05$, $\omega = 1$: geen
aantoonbare beperking --- een \omterm{def:b3:molecular-evolution:duplication}{pseudogen}, of een gen waarin \omterm{def:b3:molecular-evolution:dnds}{positieve
selectie} op sommige plaatsen de \omterm{def:b3:molecular-evolution:dnds}{zuiverende selectie} op andere in
evenwicht houdt.
\end{solution}

\begin{solution}{exo:b3:molecular-evolution:8}
$k = d/2T = 0.012/(1.3\times 10^{7}) = 9.2\times 10^{-10}$ per plaats
per jaar; $\mu = 25k = 2.3\times 10^{-8}$ per plaats per generatie;
$6\times
10^{9}\times 2.3\times 10^{-8} \approx 140$ nieuwe mutaties per diploïd
genoom per generatie. Rechtstreeks sequencen van ouders en kinderen
vindt er ongeveer zeventig: de schatting uit de divergentie is
opgeblazen door het polymorfisme dat al in de voorouderlijke populatie
aanwezig was, wat een paar honderdduizend generaties aan de schijnbare
splitsing toevoegt.
\end{solution}

\begin{solution}{exo:b3:molecular-evolution:9}
$\mathrm{e}^{-T/2N} = \mathrm{e}^{-1} = 0.37$; strijdig deel $\tfrac{2}{3}
\times 0.37 = 0.25$. Voor \qty{10}{\%}: $\mathrm{e}^{-T/2N} = 0.15$, $T =
2N\ln(1/0.15) = \num{190000}$ generaties.
\end{solution}

\begin{solution}{exo:b3:molecular-evolution:10}
Voor $s \to 0$: $1 - \mathrm{e}^{-2s} \approx 2s$ en $1 - \mathrm{e}^{-4Ns}
\approx 4Ns$, verhouding $1/2N$. Voor $4Ns \gg 1$ is de noemer $1$ en de
teller $\approx 2s$. Voor $s < 0$ met $4N|s| \gg 1$: teller $1 -
\mathrm{e}^{2|s|} \approx -2|s|$, noemer $1 - \mathrm{e}^{4N|s|} \approx
-\mathrm{e}^{4N|s|}$, dus $P \approx 2|s|\,\mathrm{e}^{-4N|s|}$.
Honderdvoudig onder neutraal bij $N = 10^{4}$: $2|s|\,\mathrm{e}^{-4\times 10^{4}|s|} =
5\times 10^{-7}$; $|s| = 1.6\times 10^{-4}$ geeft $3.2\times 10^{-4}\times
\mathrm{e}^{-6.4} = 5.3\times 10^{-7}$. Dus $|s| \approx 1.6\times 10^{-4}$,
$4N|s| \approx 6.5$: een nadeel van een honderdste procent volstaat, in
tienduizend individuen, om de fixatie honderd maal zeldzamer te maken
dan het toeval.
\end{solution}

\begin{solution}{exo:b3:molecular-evolution:11}
Met het splitsingspunt op afstand $a$ van $O$ op beide paden geldt $a + m =
0.30$, $a + h = 0.24$, $m + h = 0.20$: $m - h = 0.06$, dus $m = 0.13$, $h
= 0.07$, verhouding $1.9$. Een strenge klok per generatie zou de
verhouding van de aantallen generaties geven, $25/0.5 = 50$. De
waargenomen $1.9$ betekent dat de muis per jaar maar het dubbele van de
menselijke verandering heeft opgehoopt ondanks vijftig maal zoveel
generaties: het mutatietempo per generatie moet bij de mens zo'n $25$
maal hoger liggen --- meer celdelingen in de kiembaan per generatie ---
zodat de twee lijnen per jaar maar een factor twee verschillen.
\end{solution}

\begin{solution}{exo:b3:molecular-evolution:12}
Substituties die een stop maken komen met $1000\times 2\times 10^{-9}\times
0.04 = 8\times 10^{-8}$ per jaar: de eerste wordt na ongeveer
\qty{12}{Myr} verwacht (leesraamverschuivingen halveren dat ruwweg). In
dat venster is een gunstige mutatie, één op de $10^{3}$ tot $10^{5}$ van
alle mutaties, veel minder waarschijnlijk dan een stopcodon, en ook
wanneer zij komt wordt zij met een kans van slechts $2s$ vastgelegd.
Dus is de kopie meestal dood voordat de selectie er een nut voor kan
vinden --- de waargenomen halveringstijd van kopieën bij gewervelden is
enkele miljoenen jaren.
\end{solution}

\begin{solution}{pb:b3:molecular-evolution:1}
\textbf{1.} $k = 0.012/(2\times 6.5\times 10^{6}) = 9.2\times 10^{-10}$ per
plaats per jaar; $\mu = 25k = 2.3\times 10^{-8}$ per plaats per generatie.
\textbf{2.} $6\times 10^{9}\times 2.3\times 10^{-8} \approx 140$ nieuwe
mutaties per diploïd genoom; \qty{1.5}{\%} daarvan, ongeveer $2$, in
coderende sequentie.
\textbf{3.} Per haploïd genoom worden er $3\times 10^{9}\times 2.3\times 10^{-8}
\approx 70$ substituties per generatie vastgelegd, tegen de $1/300$ van
Haldane: twintigduizend maal meer dan de selectie zou kunnen drijven,
dus is de overgrote meerderheid neutraal.
\textbf{4.} $2N\mu = 10^{5}\times 2.3\times 10^{-8} = 2.3\times 10^{-3}$
nieuwe mutaties per plaats per generatie in de populatie; $1/2N =
10^{-5}$ ervan wordt vastgelegd.
\textbf{5.} $4N = \num{200000}$ generaties, \qty{5}{Myr} --- bijna even
lang als de tijd sinds de splitsing.
\textbf{6.} De divergentie telt de mutaties die langs de twee
afzonderlijke lijnen sinds hun gemeenschappelijke voorouder zijn
ontstaan; elk wordt binnen zijn eigen lijn vastgelegd, en het tempo op
den duur is $\mu$ hoeveel tijd elk daar ook over doet, omdat mutaties
gestaag ontstaan en de vertraging de stroom alleen uitstelt en niet
verandert. (De enige correctie geldt het polymorfisme dat er bij de
splitsing was, waardoor de gemeenschappelijke voorouder van twee
allelen ouder is dan de splitsing van de soorten.)
\textbf{7.} Neutraal $10^{-5}$; $s = 0.001$: $0.002$; $s = 0.01$: $0.0198$;
$s = 0.1$: $1 - \mathrm{e}^{-0.2} = 0.18$.
\textbf{8.} $s = 1/4N = 5\times 10^{-6}$: daaronder overspoelt de drift
de selectie en gedraagt de mutatie zich als neutraal; daarboven bepaalt
de selectie de kansen.
\textbf{9.} $s = -0.001$: $P = 0.002\,\mathrm{e}^{-200}$, voor alle
doeleinden nul --- $10^{-85}$ van neutraal. $s = -10^{-5}$, $4N|s| = 2$: $P =
2\times 10^{-5}/(\mathrm{e}^{2} - 1) = 3.1\times 10^{-6}$, \qty{31}{\%} van
neutraal: licht schadelijke mutaties worden wel degelijk vastgelegd.
\textbf{10.} Gunstige substituties per plaats per generatie: $2N\times
10^{-5}\mu\times 2s = 10^{5}\times 10^{-5}\times 0.02\,\mu = 0.02\mu$;
neutraal: $\approx\mu$. Adaptief deel $0.02/1.02 \approx \qty{2}{\%}$.
Hoewel elke gunstige mutatie tweeduizend maal waarschijnlijker wordt
vastgelegd dan een neutrale, zijn zij zo zeldzaam dat maar één op de
vijftig substituties adaptief is: de selectie werkt en is nauwelijks te
zien.
\textbf{11.} $p_{N} = 46/920 = 0.050$, $p_{S} = 84/280 = 0.30$.
Gecorrigeerd: $d_{N} = -\tfrac{3}{4}\ln(1 - 0.0667) = 0.052$, $d_{S} = -\tfrac{3}{4}
\ln(1 - 0.40) = 0.38$. $\omega = 0.14$.
\textbf{12.} \omterm{def:b3:molecular-evolution:dnds}{Zuiverende selectie} die ongeveer \qty{86}{\%} van de
veranderingen van aminozuren verwijdert --- een gewoon gen. $k = d_{S}/2T = 0.38/(1.8\times
10^{8}) = 2.1\times 10^{-9}$ per plaats per jaar: het dubbele van de
waarde voor mens en chimpansee, zoals te verwachten bij een
vergelijking die de snel tikkende lijn van het knaagdier omvat;
dezelfde orde van grootte.
\textbf{13.} $d = -\ln(1 - 0.55) = 0.80$ per plaats; $T = 0.80/(2\times
10^{-9}) = \qty{400}{Myr}$.
\textbf{14.} Twee verschillende ketens vormen de tetrameer
$\alpha_{2}\beta_{2}$, wier coöperatieve binding van zuurstof, effect
van Bohr en allosterische regeling door
2,3-bisfosfoglyceraat afhangen van het raakvlak tussen ongelijke
subeenheden; en het $\beta$-cluster kon zich daarna uitsplitsen in
embryonale, foetale en volwassen ketens met verschillende affiniteiten.
\textbf{15.} $2kT = 2\times 8\times 10^{-9}\times 4\times 10^{7} = 0.64$
substituties per plaats; met verzadiging onder twintig aminozuren is
het ruwe verschil ongeveer $0.95(1 - \mathrm{e}^{-0.64/0.95}) \approx 0.47$.
Bij \qty{500}{Myr} is $d = 8$: elke plaats is vele malen veranderd en
zit het ruwe verschil aan zijn plafond van \qty{95}{\%}, zonder nog
iets over de tijd te zeggen.
\textbf{16.} $10^{-11}\times 10^{9}\times 100\times 2 = 2$ substituties
tussen twee lijnen in 100 plaatsen over een miljard jaar. Voor een
splitsing van \qty{10}{Myr} is de verwachting $0.02$: vrijwel altijd
geen enkel verschil, dus geen oplossend vermogen.
\textbf{17.} $1200\times 2\times 10^{-9}\times 0.04 = 9.6\times 10^{-8}$
per jaar: de eerste stop wordt na ongeveer \qty{10}{Myr} verwacht.
\textbf{18.} Het hele genoom verdubbelen verdubbelt elk gen tegelijk,
zodat de stoichiometrie van op elkaar inwerkende eiwitten behouden
blijft en geen onevenwicht in de dosis tegen de kopieën selecteert; een
losse tandemkopie verandert de dosis van één gen en is vaak schadelijk,
of onmiddellijk overtollig en vrij om te vervallen.
\textbf{19.} $\mathrm{e}^{-\num{80000}/\num{100000}} = \mathrm{e}^{-0.8} = 0.45$.
\textbf{20.} $\tfrac{2}{3}\times 0.45 = 0.30$: in overeenstemming met de
waargenomen \qty{30}{\%}.
\textbf{21.} Elk de helft: \qty{15}{\%} van alle bomen groepeert de mens
met de gorilla en \qty{15}{\%} de chimpansee met de gorilla. Een
duidelijke overmaat aan één van beide zou wijzen op genstroom tussen die
twee soorten na hun splitsing, die de drift niet kan voortbrengen.
\textbf{22.} $N = \num{10000}$: $\mathrm{e}^{-4} = 0.018$, strijdigheid
\qty{1.2}{\%}. De waargenomen \qty{30}{\%} vergt dus een voorouderlijke
populatie van zo'n vijftigduizend --- enkele malen de effectieve omvang
van de mens vandaag.
\textbf{23.} Aaneenschakelen veronderstelt dat elk gen de \omterm{prop:b3:molecular-evolution:ils}{soortboom}
heeft; bij \omterm{prop:b3:molecular-evolution:ils}{onvolledige sortering van lijnen} hebben de genen vele bomen,
en bij korte inwendige takken kan de meest voorkomende \omterm{prop:b3:molecular-evolution:ils}{genboom} van de
\omterm{prop:b3:molecular-evolution:ils}{soortboom} afwijken, zodat meer gegevens het verkeerde antwoord
zelfverzekerder maken. Een methode met de coalescent berekent voor elke
\omterm{prop:b3:molecular-evolution:ils}{kandidaat-soortboom} de verwachte verdeling van \omterm{prop:b3:molecular-evolution:ils}{genbomen}, en kiest de
\omterm{prop:b3:molecular-evolution:ils}{soortboom} die de waargenomen frequenties het best verklaart.
\textbf{24.} De \omterm{prop:b3:molecular-evolution:ils}{onvolledige sortering van lijnen} gebeurde in de
gemeenschappelijke voorouderlijke populatie van alle moderne mensen, dus
zou zij elke menselijke populatie even verwant aan de neanderthalers
maken. Een overmaat bij niet-Afrikanen betekent dat neanderthaler-allelen
zijn binnengekomen nadat de voorouders van de niet-Afrikanen Afrika
hadden verlaten: vermenging, zo'n vijftigduizend jaar geleden.
\textbf{25.} $\mu \approx 2.3\times 10^{-8}$ per plaats per generatie;
$\omega \approx 0.14$; ongeveer \qty{30}{\%} van de \omterm{prop:b3:molecular-evolution:ils}{genbomen} is met de
\omterm{prop:b3:molecular-evolution:ils}{soortboom} in strijd.
\end{solution}
